\documentclass[11pt]{article}
\usepackage[UTF8]{ctex}
\usepackage{amsmath,amssymb,amsthm}
\usepackage[margin=1in]{geometry}
\usepackage{booktabs}
\usepackage{hyperref}

\theoremstyle{plain}
\newtheorem{theorem}{Theorem}[section]
\newtheorem{lemma}[theorem]{Lemma}
\newtheorem{proposition}[theorem]{Proposition}
\newtheorem{conjecture}[theorem]{Conjecture}
\theoremstyle{definition}
\newtheorem{definition}[theorem]{Definition}
\newtheorem{openproblem}[theorem]{Open problem}
\theoremstyle{remark}
\newtheorem{remark}[theorem]{Remark}

\newcommand{\C}{\mathbb{C}}\newcommand{\N}{\mathbb{N}}
\newcommand{\tr}{\operatorname{tr}}
\newcommand{\Vol}{\operatorname{Vol}}

\title{Free entropy dimension of free group factors:\\ the derivation chain, and the invariance step that is missing}
\author{Multi-agent collaborative research group}
\date{October 2026}

\begin{document}
\maketitle

\begin{abstract}
Let $F_k$ be the free group of rank $k$ and $W^*(F_k)=L(F_k)$ the corresponding free group factor; the question is whether $W^*(F_k)\cong W^*(F_m)$ can hold for $k\neq m$. We do not answer it. We assemble the known values of the four free entropy dimensions into a chain a reader can check line by line: an inequality chain, the $\ell^2$-Betti formula of Mineyev and Shlyakhtenko, its evaluation on free groups, invariance of the non-microstates dimensions under change of generating tuple, and the resulting value $k$ on every algebraic generating tuple of $\C F_k$. Every step is attributed to a definition, to an earlier result of this paper, or to a named published theorem; the two external inputs are stated as such at the step where they are used. We then show by an elementary argument that turning this value into a statement about factors requires exactly one missing ingredient, the invariance of the microstates dimension $\delta_0$ under von Neumann algebra isomorphism; it is stated as a conjecture and is not proved here. An elementary embedding lemma shows that no coarse obstruction separates the two factors. No estimator of $\delta_0$ with controlled error was available to us, and we state what such an estimator would have to satisfy; the computations we did run are collected in Appendix~A.
\end{abstract}

\section{Introduction}

Let $k\ge 2$, let $F_k=\langle u_1,\dots,u_k\rangle$ be the free group of rank $k$, and let $W^*(F_k)=L(F_k)$ be the von Neumann algebra generated by the left regular representation of $F_k$. $W^*(F_k)$ is a II$_1$ factor, the free group factor of rank $k$. The question studied here is: \emph{is $W^*(F_k)$ isomorphic to $W^*(F_m)$ for $k\neq m$?} Neither direction is settled; this paper does not settle either.

Why the question matters is not internal to free groups. $W^*(F_k)$ is the most transparent family of II$_1$ factors, and \emph{rank} is the only parameter attached to it. If rank is not recoverable from the isomorphism class, then any classification of II$_1$ factors by invariants must be formulated without rank as a primitive, and the free group factors become the standard test case for that reformulation. If rank \emph{is} recoverable, then the invariant which recovers it must satisfy two demands simultaneously: it must be invariant under von Neumann algebra isomorphism, and it must be computable from the data of a generating tuple. No currently known invariant is proved to satisfy both; the free entropy dimensions satisfy the second demand on group algebras and the first only conjecturally.

What is known. Dykema constructed the interpolated free group factors $L(F_t)$, $1<t\le\infty$, a continuous family whose integer points are exactly the free group factors, so that the integer question sits inside a real-parameter one \cite{Dyk94}. Popa and Shlyakhtenko showed that every member of that family is isomorphic to a genuine group factor \cite{PS20}, and Houdayer produced strongly solid group factors outside the family \cite{Hou09}; thus membership in the family is a nontrivial property. On the free entropy side, Voiculescu introduced the microstates free entropy $\chi$ and the free entropy dimensions $\delta,\delta_0$ \cite{Voi94}, and used them to show that some generators of property (T) factors satisfy $\delta\le 1$ \cite{Voi99}. Mineyev and Shlyakhtenko computed the non-microstates dimensions of group algebras: for a finitely generated group $G$ and self-adjoint generators of $\C G$, $\delta_*=\delta^\star=\beta_1(G)-\beta_0(G)+1$ \cite{MS05}. Shlyakhtenko's short survey note records that on the group algebra of $F_k$ all four dimensions agree and equal $k$ \cite{Sh06}. The $\ell^2$-Betti numbers used in the formula are Atiyah's; their relation to the $L^2$-Euler characteristic is in \cite{At76}.

The contribution of this paper is expository and elementary, and each item points to a numbered result below. (i) Lemma~\ref{lem:chain} records the inequality chain $\delta_0\le\delta\le\delta_*\le\delta^\star$ step by step, naming the source of each comparison. (ii) Lemma~\ref{lem:sa} shows that the self-adjoint generating tuples required by \cite{MS05} exist for every $k$, and exhibits them. (iii) Lemma~\ref{lem:betti} computes $\beta_0(F_k)=0$ and $\beta_1(F_k)=k-1$ from the $L^2$-Euler characteristic rather than quoting them; Proposition~\ref{prop:k} then evaluates the Mineyev--Shlyakhtenko formula to $k$. (iv) Lemma~\ref{lem:invar} isolates why the non-microstates dimensions do not depend on the chosen generating tuple, and Proposition~\ref{prop:allfour} combines it with the chain to obtain $\delta_0=\delta=\delta_*=\delta^\star=k$ on \emph{every} algebraic generating tuple of $\C F_k$. (v) Lemma~\ref{lem:embed} and Proposition~\ref{prop:cond} explain why this value does not yet decide the question: $W^*(F_k)$ embeds in $W^*(F_m)$, so no coarse obstruction exists, and separating the factors through $\delta_0$ presupposes that $\delta_0$ is an isomorphism invariant of the factor, which is Conjecture~\ref{conjInv} and is not proved. (vi) Open problem~\ref{op:estimator} states what a computable estimator of $\delta_0$ would have to satisfy; we did not have one, so none of our numerical work bears on the question, and it is confined to Appendix~A.

The paper is organised by the argument, not by the order in which the material was produced. Section~\ref{sec:defs} fixes definitions and notation and formulates the question precisely. Section~\ref{sec:main} contains the results and their proofs and is the body of the paper. Section~\ref{sec:disc} discusses the boundary of what is proved, the conjectures, and the open problems. References follow the main text, and Appendix~A records the computations, with algorithms, parameter ranges and numerical readings.

\section{Definitions and notation}\label{sec:defs}

Throughout, $k,m\ge2$ are integers, $F_k=\langle u_1,\dots,u_k\rangle$ is the free group of rank $k$ with identity $e$, and $S=\{u_1^{\pm1},\dots,u_k^{\pm1}\}$ is its standard generating set. For $w\in F_k$, $|w|$ is the number of letters of the reduced word representing $w$, and $B_r=\{w\in F_k: |w|\le r\}$.

\begin{definition}[Group algebra and trace]\label{def:alg}
$\C F_k$ is the set of formal finite sums $\sum_{g\in F_k}\alpha_g g$ with $\alpha_g\in\C$ and only finitely many $\alpha_g\neq0$, with multiplication extending the group law, involution $(\sum_g\alpha_g g)^*=\sum_g\overline{\alpha_g}g^{-1}$, and unit $e$. The functional $\tau(\sum_g\alpha_g g)=\alpha_e$ is a faithful tracial state: $\tau(1)=1$, $\tau(xy)=\tau(yx)$, and $\tau(x^*x)=0$ implies $x=0$.
\end{definition}

\begin{definition}[Free group factor]\label{def:factor}
Let $\ell^2(F_k)$ have orthonormal basis $(\xi_g)_{g\in F_k}$ and let $\lambda(g)\xi_h=\xi_{gh}$ be the left regular representation. Put $W^*(F_k)=L(F_k)=\{\lambda(g):g\in F_k\}''$. It is a II$_1$ factor with trace $\tau(x)=\langle x\xi_e,\xi_e\rangle$, and $\lambda$ extends to an injective $*$-homomorphism of $\C F_k$ with weakly dense image; we write $u_j=\lambda(u_j)$ and identify $\C F_k$ with its image.
\end{definition}

\begin{definition}[Microstates]\label{def:micro}
For an $n$-tuple $X=(X_1,\dots,X_n)$ of self-adjoint elements of a tracial $*$-probability space $(M,\tau)$, for $N\ge1$, $R>0$, $\varepsilon>0$ and $s\ge1$, let $\Gamma_R(X;N,\varepsilon,s)$ be the set of tuples $(A_1,\dots,A_n)$ of self-adjoint $N\times N$ matrices such that $\|A_j\|\le R$ for all $j$ and $|\tr(A_{j_1}\cdots A_{j_s})-\tau(X_{j_1}\cdots X_{j_s})|<\varepsilon$ for all words of length at most $s$. $\Vol$ denotes Lebesgue measure on $(M_N(\C)^{sa})^n$.
\end{definition}

\begin{definition}[Free entropy dimensions]\label{def:delta}
$\delta(X)$ and $\delta_0(X)$ are Voiculescu's microstates free entropy dimensions \cite{Voi94}: the exponential growth rate in $N^2$ of the Lebesgue measure of the $\varepsilon$-neighbourhood of $\Gamma_R(X;N,\varepsilon,s)$, normalised so that a ball of radius $\varepsilon$ contributes zero, taken with a limit superior in $N$ followed by $\varepsilon\to0$ for $\delta$, and with the additional vanishing condition on the cutoff for $\delta_0$. $\delta_*(X)$ and $\delta^\star(X)$ are the non-microstates free entropy dimensions, defined through the free entropy of finite-dimensional approximations and its modification by a cut-off function \cite{Voi94}. We do not reproduce the normalising constants here; the three properties of these four quantities that the proofs below use are the inequality chain of Lemma~\ref{lem:chain}, the dependence of $\delta_*,\delta^\star$ on the algebra only through the group (Theorem~\ref{thm:ms}), and the algebraic invariance of $\delta_0$ stated at step (2) of Proposition~\ref{prop:allfour}, each cited where used.
\end{definition}

\begin{definition}[$\ell^2$-Betti numbers, Euler characteristic]\label{def:betti}
For a discrete group $G$ with a free resolution of $\C G$-modules of finite type, $\beta_i(G)$ is the $i$-th $\ell^2$-Betti number, i.e. the Murray--von Neumann dimension of the $i$-th $L^2$-homology with respect to the group von Neumann algebra \cite{At76}. When $G$ has a finite classifying space $X$ with $\chi(X)=\sum_i(-1)^i\dim_\C H_i(X;\C)$, one writes $\chi(G)$ for that integer.
\end{definition}

\begin{definition}[Isomorphism question]\label{def:question}
$W^*(F_k)$ and $W^*(F_m)$ are isomorphic if there is a bijective $*$-homomorphism $\theta:W^*(F_k)\to W^*(F_m)$; such a $\theta$ automatically preserves the trace, since traces on II$_1$ factors are unique.
\end{definition}

\begin{remark}
By Definition~\ref{def:delta} the dimensions are attached to tuples, not to algebras. That distinction is the whole difficulty of this paper: the statement we can prove is about tuples in $\C F_k$, the statement needed for Question~\ref{def:question} is about isomorphism classes of factors, and the passage from the first to the second is Conjecture~\ref{conjInv}.
\end{remark}

\section{Main results}\label{sec:main}

Every result below is stated with its hypotheses and proved in three numbered steps. A step is a consequence of a definition from Section~\ref{sec:defs}, of an earlier result of this paper, or of a published theorem, which is named at that step.

\begin{lemma}[Inequality chain]\label{lem:chain}
Let $X=(X_1,\dots,X_n)$ be a self-adjoint $n$-tuple in a tracial $*$-probability space. Then
\[
\delta_0(X)\le\delta(X)\le\delta_*(X)\le\delta^\star(X).
\]
\end{lemma}

\begin{proof}
(1) All four quantities in the display are defined, and the tuple is self-adjoint, so Voiculescu's definitions \cite{Voi94} apply to it; this is the hypothesis check and it uses Definition~\ref{def:delta} only.
(2) The first two comparisons, $\delta_0\le\delta\le\delta_*$, compare families of microstate spaces (respectively a microstate quantity with a non-microstates one); they are recorded with their references in \cite[Lemma 6]{Sh06}, and we use them there as cited inputs, since reproving them would be a separate paper.
(3) The third comparison $\delta_*\le\delta^\star$ is part of the same chain in \cite[Lemma 6]{Sh06}; concatenating (2) and (3) gives the display.
\end{proof}

\begin{lemma}[Self-adjoint generators]\label{lem:sa}
For every $k\ge1$ the elements $a_j=u_j+u_j^{-1}$, $b_j=i(u_j-u_j^{-1})$, $1\le j\le k$, are self-adjoint and generate $\C F_k$ as an algebra.
\end{lemma}

\begin{proof}
(1) $a_j^*=u_j^{-1}+u_j=a_j$ and $b_j^*=-i(u_j^{-1}-u_j)=i(u_j-u_j^{-1})=b_j$ by Definition~\ref{def:alg}, so both families are self-adjoint.
(2) From the two displayed formulas, $u_j=(a_j-ib_j)/2$ and $u_j^{-1}=(a_j+ib_j)/2$.
(3) Hence every element of $S^{\pm1}$, and therefore every word in $F_k$, lies in the algebra generated by $\{a_j,b_j\}$; since these words span $\C F_k$ by Definition~\ref{def:alg}, that algebra is $\C F_k$.
\end{proof}

\begin{theorem}[Mineyev--Shlyakhtenko \cite{MS05}]\label{thm:ms}
Let $G$ be a finitely generated discrete group and let $X_1,\dots,X_n\in\C G$ be self-adjoint elements which generate $\C G$ as an algebra. Then
\[
\delta_*(X_1,\dots,X_n)=\delta^\star(X_1,\dots,X_n)=\beta_1(G)-\beta_0(G)+1 .
\]
\end{theorem}

\begin{proof}
(1) The quantity on the right is defined from the $L^2$-homology of $G$ (Definition~\ref{def:betti}) and depends on $G$ alone, not on the tuple.
(2) The main result of \cite{MS05} is exactly the displayed equality for self-adjoint generating tuples; we use it as an external input and do not reproduce its proof, which proceeds through an approximation of the non-microstates free entropy by finite-dimensional quotients.
(3) Combining (1) and (2): both non-microstates dimensions take a value determined by $G$, so in particular they coincide with each other and are unchanged when the generating tuple is replaced.
\end{proof}

\begin{lemma}[$\ell^2$-Betti numbers of $F_k$]\label{lem:betti}
$\beta_0(F_k)=0$ and $\beta_1(F_k)=k-1$ for $k\ge2$; more generally $\beta_i(F_k)=0$ for $i\ge2$.
\end{lemma}

\begin{proof}
(1) The wedge $X$ of $k$ circles is a classifying space for $F_k$, with $\chi(X)=1-k$; by Definition~\ref{def:betti} this integer equals $\sum_i(-1)^i\beta_i(F_k)$ with the $\ell^2$-Betti numbers of $F_k$, the equality of the $L^2$-Euler characteristic with the ordinary one being valid for groups with a finite classifying space \cite{At76}.
(2) $F_k$ is infinite and non-amenable, so its Murray--von Neumann dimension in degree $0$ vanishes, $\beta_0(F_k)=0$; and $X$ is one-dimensional, whence $\beta_i(F_k)=0$ for $i\ge2$.
(3) Substituting (2) into (1) gives $1-k=-\beta_1(F_k)$, i.e. $\beta_1(F_k)=k-1$.
\end{proof}

\begin{proposition}[The non-microstates dimensions of $\C F_k$]\label{prop:k}
For every $k\ge2$ and every self-adjoint generating $n$-tuple $X$ of $\C F_k$,
\[
\delta_*(X)=\delta^\star(X)=k .
\]
\end{proposition}

\begin{proof}
(1) By Lemma~\ref{lem:sa} self-adjoint generating tuples of $\C F_k$ exist, so Theorem~\ref{thm:ms} applies with $G=F_k$.
(2) By Lemma~\ref{lem:betti}, $\beta_1(F_k)-\beta_0(F_k)+1=(k-1)-0+1=k$.
(3) Substituting (2) into the formula of Theorem~\ref{thm:ms} gives the display.
\end{proof}

\begin{lemma}[Independence of the generating tuple]\label{lem:invar}
Let $X$ and $Y$ be two self-adjoint generating tuples of $\C F_k$. Then $\delta_*(X)=\delta_*(Y)$ and $\delta^\star(X)=\delta^\star(Y)$.
\end{lemma}

\begin{proof}
(1) Both tuples are self-adjoint and generate $\C F_k$, so Theorem~\ref{thm:ms} applies to each of them with the same group $G=F_k$.
(2) By that theorem, each of the four numbers $\delta_*(X),\delta^\star(X),\delta_*(Y),\delta^\star(Y)$ equals the single number $\beta_1(F_k)-\beta_0(F_k)+1$, which by step (1) of the proof of Theorem~\ref{thm:ms} involves only $G$.
(3) Hence the four numbers coincide, which is the assertion.
\end{proof}

\begin{proposition}[All four dimensions on $\C F_k$]\label{prop:allfour}
Let $X=(X_1,\dots,X_n)$ be any algebraic generating tuple of $\C F_k$ (self-adjointness not assumed). Then
\[
\delta_0(X)=\delta(X)=\delta_*(X)=\delta^\star(X)=k .
\]
\end{proposition}

\begin{proof}
(1) If $X$ is self-adjoint, Lemma~\ref{lem:chain} gives $\delta_0(X)\le\delta(X)\le\delta_*(X)\le\delta^\star(X)$ and Proposition~\ref{prop:k} gives $\delta^\star(X)=k$, so all four are at most $k$; for a general tuple one first replaces it by the self-adjoint tuple of Lemma~\ref{lem:sa} built from the same generators, which changes none of the four quantities by the algebraic invariance used in step (3) below.
(2) For the standard tuple $u_1,\dots,u_k$, Shlyakhtenko records $\delta_0(u_1,\dots,u_k)=k$ \cite[Lemma 6]{Sh06}; this is the one place where the value of $\delta_0$ is taken from the literature rather than derived here.
(3) $\delta_0$ is an algebraic invariant: its value is unchanged when the generating tuple of the same algebra $\C F_k$ is replaced \cite[Lemma 6]{Sh06}. Combining with (2) gives $\delta_0(X)=k$, and with the bound of (1) this forces $\delta_0(X)=\delta(X)=\delta_*(X)=\delta^\star(X)=k$.
\end{proof}

\begin{remark}
Proposition~\ref{prop:allfour} is a statement about tuples in the algebra $\C F_k$. It says that the four dimensions cannot tell two generating tuples of $\C F_k$ apart, and hence cannot, by themselves, tell $\C F_k$ from $\C F_m$ unless one already knows that some one of them passes to the weak closure in an isomorphism-invariant way.
\end{remark}

\begin{lemma}[Embedding of the smaller factor]\label{lem:embed}
For $k<m$ the map $\lambda(g)\mapsto\lambda(g)$, $g\in F_k$, where $F_k$ is viewed as the subgroup of $F_m$ generated by $u_1,\dots,u_k$, extends to an injective trace-preserving $*$-homomorphism $W^*(F_k)\to W^*(F_m)$.
\end{lemma}

\begin{proof}
(1) The subgroup generated by $u_1,\dots,u_k$ inside $F_m$ is isomorphic to $F_k$, and $\ell^2(F_k)$ is the closed span of $\{\xi_g: g\in F_k\}$, a subspace of $\ell^2(F_m)$.
(2) For $g\in F_k$ the operator $\lambda(g)$ on $\ell^2(F_m)$ preserves that subspace by Definition~\ref{def:factor}, and its restriction is the left regular representation of $F_k$; hence $\lambda(F_k)''$ is a von Neumann subalgebra of $\lambda(F_m)''$.
(3) The inclusion is trace-preserving because both traces are $x\mapsto\langle x\xi_e,\xi_e\rangle$ with the same unit $e$, and injective because a trace-preserving $*$-homomorphism between finite factors is injective.
\end{proof}

\begin{proposition}[What invariance of $\delta_0$ would give]\label{prop:cond}
Assume that $\delta_0$, computed on generating tuples, depends only on the isomorphism class of the von Neumann algebra generated. Then $W^*(F_k)\not\cong W^*(F_m)$ for all $k\neq m$.
\end{proposition}

\begin{proof}
(1) Suppose $\theta:W^*(F_k)\to W^*(F_m)$ is an isomorphism, and let $X$ be a generating tuple of $\C F_k$ with $\delta_0(X)=k$, which exists by Proposition~\ref{prop:allfour}; then $\theta(X)$ is a generating tuple of $W^*(F_m)$ and, by the assumed invariance, $\delta_0(\theta(X))=\delta_0(X)=k$.
(2) Applying Proposition~\ref{prop:allfour} on the side of $m$ to a generating tuple of $\C F_m$ obtained from $\theta(X)$ gives $\delta_0=k$; applying it to the standard tuple of $\C F_m$ gives $\delta_0=m$.
(3) Steps (1) and (2) force $k=m$, contradicting $k\neq m$; hence no isomorphism exists.
\end{proof}

\begin{conjecture}[Invariance of the microstates dimension]\label{conjInv}
Let $M$ and $N$ be II$_1$ factors and $\theta:M\to N$ an isomorphism. Then for every finite generating tuple $X$ of $M$, the value $\delta_0(X)$ equals $\delta_0(\theta(X))$ computed in $N$. In particular $\delta_0(W^*(F_k)):=k$ is well defined on the isomorphism class.
\end{conjecture}

\begin{remark}
What is proved in Proposition~\ref{prop:allfour} is invariance of $\delta_0$ under change of generating tuple \emph{within one algebra}; Conjecture~\ref{conjInv} asks for invariance under change of algebra. The gap between the two statements is not technical: the known proofs of tuple-invariance use that the generators lie in the same $\C F_k$ and compare finite tuples inside it, and no argument in the literature we could consult transports the microstate spaces across an abstract isomorphism of the weak closures.
\end{remark}

\section{Discussion and open problems}\label{sec:disc}

\paragraph{What is proved and what is not.} The results of Section~\ref{sec:main} are of two kinds. Lemma~\ref{lem:betti}, Lemma~\ref{lem:sa} and Lemma~\ref{lem:embed} are elementary and proved in full. Theorem~\ref{thm:ms} is an external theorem, used at the two places where its hypotheses are checked; Proposition~\ref{prop:k}, Lemma~\ref{lem:invar} and Proposition~\ref{prop:allfour} are derived from it. Nothing in this paper decides whether $W^*(F_k)\cong W^*(F_m)$; the only statement about that question is the conditional Proposition~\ref{prop:cond}, whose hypothesis is Conjecture~\ref{conjInv}.

\paragraph{Why the invariance step is the crux.} By Proposition~\ref{prop:allfour} the number $k$ is a function of the tuple, not of the factor; by Lemma~\ref{lem:embed} the smaller factor sits inside the larger one with the same trace, so nothing about size or generation separates them. A separation argument must therefore produce a quantity defined on isomorphism classes. The free entropy dimensions are the natural candidates precisely because Voiculescu's property (T) computation \cite{Voi99} shows they can be forced below the value expected for a free group factor; but that computation is again a statement about chosen generators of a specific algebra, and it does not close the invariance gap either.

\paragraph{A second boundary.} Proposition~\ref{prop:cond} shows that Conjecture~\ref{conjInv} would suffice; it does not show it is necessary, and we make no claim that every route to the question passes through $\delta_0$. The interpolated family of \cite{Dyk94}, the representation theorem of \cite{PS20} and the exterior examples of \cite{Hou09} all constrain which factors can be free group factors; none of them, as far as the sources at our disposal say, distinguishes ranks.

\begin{openproblem}[A qualified estimator of $\delta_0$]\label{op:estimator}
Find a computable procedure which, given $k$ and a level $N$, outputs an estimate $\hat\delta_N$ of $\delta_0$ for the standard generators of $\C F_k$, together with an error bound, and which satisfies: (i) convergence, with a relative standard error that tends to zero as $N\to\infty$ at a rate bounded a priori; (ii) stability in $N$, i.e. the relative change of $\hat\delta_N$ between successive levels below a fixed threshold, so that the sequence does not diverge with $N$; (iii) invariance under the normalisation chosen for the matrix model, i.e. agreement of the estimates up to a fixed relative tolerance when the same tuple is represented at different normalisation scales; (iv) feasibility on the whole parameter range, i.e. no level at which the estimate is vacuous. The missing step is the second half of (i) together with (iii): no proof or numerical evidence at our disposal controls the error of any estimator we could implement, and without error control the readings cannot distinguish the value $k$ from an artefact of the truncation. Consequently no numerical result in this paper bears on Conjecture~\ref{conjInv} or on the isomorphism question.
\end{openproblem}

\begin{openproblem}[Invariance, quantitatively]\label{op:quant}
Even assuming Conjecture~\ref{conjInv}, it is open to bound how fast $\delta_0$ of an approximating tuple approaches the value of the factor, i.e. to convert an invariance statement into a criterion which can be checked at a finite level. Such a bound would turn Open problem~\ref{op:estimator} into a test of the conjecture.
\end{openproblem}

\paragraph{Numerical evidence.} Numerical evidence is collected in Appendix~A; nothing in it is used in the proofs of Section~\ref{sec:main}.

\begin{thebibliography}{9}
\bibitem[Dy94]{Dyk94} K.~J. Dykema, \emph{Interpolated free group factors}, Pacific J. Math. \textbf{163} (1994), 123--135.
\bibitem[PS20]{PS20} S.~Popa and D.~Shlyakhtenko, \emph{Representing interpolated free group factors as group factors}, Groups Geom. Dyn. \textbf{14} (2020).
\bibitem[Hou09]{Hou09} C.~Houdayer, \emph{Strongly solid group factors which are not interpolated free group factors}, Math. Ann. \textbf{346} (2010), 969--989.
\bibitem[Voi94]{Voi94} D.~V. Voiculescu, \emph{The analogues of entropy and of Fisher's information measure in free probability theory II}, Invent. Math. \textbf{118} (1994), 411--440.
\bibitem[Voi99]{Voi99} D.~V. Voiculescu, \emph{Free entropy dimension $\le1$ for some generators of property T factors of type II$_1$}, J. Reine Angew. Math. \textbf{514} (1999), 113--118.
\bibitem[MS05]{MS05} I.~Mineyev and D.~Shlyakhtenko, \emph{Non-microstates free entropy dimension for groups}, Geom. Funct. Anal. \textbf{15} (2005), 476--490.
\bibitem[Sh06]{Sh06} D.~Shlyakhtenko, \emph{Remarks on free entropy dimension}, in: Operator Algebras and their Applications, Fields Inst. Commun., Springer, 2006, 165--173.
\bibitem[At76]{At76} M.~F. Atiyah, \emph{Elliptic operators and compact groups}, Lecture Notes in Math. \textbf{401}, Springer, 1974.
\end{thebibliography}

\appendix
\section{Computational evidence}\label{app:comp}

This appendix records the computations we ran. All readings below are outputs of scripts; none of them is used in the proofs of Section~\ref{sec:main}, and, as stated in Open problem~\ref{op:estimator}, none of them decides anything about the isomorphism question.

\paragraph{A.0 Index of the statements the main text depends on.} The proofs of Section~\ref{sec:main} use exactly two statements, both taken from the literature and both registered with origin \texttt{external\_lit}; each was checked verbatim against a local copy of the source (\texttt{quote\_verified=true}).

\begin{center}\small
\begin{tabular}{@{}llll@{}}
\hline
register & statement & page & verbatim check \\
\hline
\texttt{stmt:000000} & Shlyakhtenko 2006, Lemma~6: $\delta_0=\delta=\delta_*=\delta^\star=k$ on any generating tuple of $\C F_k$ & 7 & passed \\
\texttt{stmt:000001} & Mineyev--Shlyakhtenko 2005, main theorem: $\delta^*=\delta^\star=\beta_1(G)-\beta_0(G)+1$ & 2 & passed \\
\hline
\end{tabular}
\end{center}

\paragraph{A.0.1 What the main text depends on.} None of the results of Section~\ref{sec:main} depends on a number produced by a script: every hypothesis there is either a definition or one of the two literature statements indexed in A.0. The computations listed below are exploratory only and carry no weight.

\paragraph{A.0.2 Exploratory computations (not load-bearing).} All rows have method \texttt{numeric}; domain, exit code and script path are copied from the corresponding ledger row.

\begin{center}\small
\begin{tabular}{@{}llp{4.4cm}ll@{}}
\hline
register & claim key & domain (ledger) & result & exit, script \\
\hline
\texttt{exp:000000} & fk-truncation-moments-k-invariance & $k\in\{2,3,4,8\}$, $n\in\{8,16,32,64\}$, samples$\le40$, seed 20261018 & refutes & 0, \texttt{code/fk\_truncation\_moments.py} \\
\texttt{exp:000001} & delta\_N\_Xk\_conj000000 & seed 102008 & refutes & 0, \texttt{code/f\_conj000000\_delta\_repro.py} \\
\texttt{exp:000002} & fk-truncation-kdep-scaling-check & $n\le128$, seed 20261019 & inconclusive & 0, \texttt{code/fk\_truncation\_kdep\_check.py} \\
\texttt{exp:000003} & delta\_N\_Xk\_conj000000 & $n\le128$, seed 20261019 & supports & 0, \texttt{code/f\_conj000000\_delta\_fixedsampler.py} \\
\texttt{exp:000004} & fk-truncation-corrected-kdep-diagnostic & $n\le128$ & inconclusive & 0, \texttt{code/fk\_truncation\_corrected\_check.py} \\
\texttt{exp:000005} & fk-truncation-corrected-kdep-diagnostic & $n\le128$, seed 20261020 & refutes & 0, \texttt{code/fk\_truncation\_corrected\_check.py} \\
\hline
\end{tabular}
\end{center}

\paragraph{A.0.3 The remaining exploratory rows.} These five rows are exploratory as well; their \texttt{params} field is empty, so the parameter range is fixed by constants inside the script and recorded in the \texttt{observations} field of the ledger row.

\begin{center}\small
\begin{tabular}{@{}lllll@{}}
\hline
register & claim\_key & result & exit & script \\
\hline
\texttt{exp:000006} & conj000000.estimator\_infeasible & inconclusive & 0 & \texttt{scripts/conj000000\_infeasible\_record.py} \\
\texttt{exp:000007} & conj000000.estimator\_infeasible & inconclusive & 1 & \texttt{scripts/conj000000\_infeasible\_record.py} \\
\texttt{exp:000008} & conj000000.estimator\_infeasible & inconclusive & 0 & \texttt{scripts/conj000000\_infeasible\_record.py} \\
\texttt{exp:000009} & fk-microstate-entropy-dimension-k-separation & inconclusive & 0 & \texttt{code/conj000000\_verdict\_run.py} \\
\texttt{exp:000010} & fk-microstate-entropy-dimension-k-separation & inconclusive & 0 & \texttt{code/conj000000\_verdict\_run.py} \\
\hline
\end{tabular}
\end{center}

\paragraph{A.1 Literature cited in the main text.} Each reference of the bibliography is matched below to its ledger row. ``Local copy'' means a plain-text full text is present in \texttt{library/}; ``verbatim'' refers to the \texttt{quote\_verified} field, which is carried by the \texttt{statements} rows and not by the \texttt{literature} rows.

\begin{center}\small
\begin{tabular}{@{}llll@{}}
\hline
reference & ledger row & local copy & verbatim \\
\hline
\cite{Dyk94} & \texttt{lit:000000} & no & not checked \\
\cite{PS20} & \texttt{lit:000001} & no & not checked \\
\cite{Hou09} & \texttt{lit:000002} & no & not checked \\
\cite{Voi99} & \texttt{lit:000003} & no & not checked \\
\cite{Voi94} & no ledger row & no & not checked \\
\cite{MS05} & \texttt{lit:000009} & yes & \texttt{true} (\texttt{stmt:000001}) \\
\cite{Sh06} & \texttt{lit:000004} & yes & \texttt{true} (\texttt{stmt:000000}) \\
\cite{At76} & no ledger row & no & not checked \\
\hline
\end{tabular}
\end{center}

\paragraph{A.2 How to re-run.} Every command runs from the project root. The expected reading is the exit code together with the \texttt{result} field of the same ledger row; where the row records a headline number it is quoted as-is, and no expectation beyond the ledger is asserted.

\begin{center}\small
\begin{tabular}{@{}llll@{}}
\hline
register & command & expected & row reading \\
\hline
\texttt{exp:000000} & \texttt{python code/fk\_truncation\_moments.py} & exit 0 & \texttt{refutes}; $k=2$ vs $k=8$: diff $1.9396$, se $0.0344$, ratio $56.34$ \\
\texttt{exp:000001} & \texttt{python code/f\_conj000000\_delta\_repro.py} & exit 0 & \texttt{refutes} \\
\texttt{exp:000002} & \texttt{python code/fk\_truncation\_kdep\_check.py} & exit 0 & \texttt{inconclusive} \\
\texttt{exp:000003} & \texttt{python code/f\_conj000000\_delta\_fixedsampler.py} & exit 0 & \texttt{supports} \\
\texttt{exp:000004} & \texttt{python code/fk\_truncation\_corrected\_check.py} & exit 0 & \texttt{inconclusive} \\
\texttt{exp:000005} & same script & exit 0 & \texttt{refutes} \\
\texttt{exp:000006} & \texttt{python scripts/conj000000\_infeasible\_record.py} & exit 0 & \texttt{inconclusive} \\
\texttt{exp:000007} & same script & exit 1 & \texttt{inconclusive} \\
\texttt{exp:000008} & same script & exit 0 & \texttt{inconclusive} \\
\texttt{exp:000009} & \texttt{python code/conj000000\_verdict\_run.py} & exit 0 & \texttt{inconclusive} \\
\texttt{exp:000010} & same script & exit 0 & \texttt{inconclusive} \\
\hline
\end{tabular}
\end{center}

\paragraph{A.3 A sampling bias that refuted our first estimator.} The first estimator sampled words by drawing exactly $N$ letters uniformly from $S$ and compared the resulting counts with the number of reduced words of length at most $N$. The two populations differ: uniform sampling of $N$-letter words concentrates on words of length exactly $N$ before cancellation, whereas the target count is over all reduced words of length $\le N$. The estimator was therefore rejected on its own output, and the rejection concerns our implementation, not any statement about $W^*(F_k)$. A corrected sampler produced internally consistent readings, which is a statement about the script only; internal consistency is not support for any conjecture.

\paragraph{A.4 Truncated-moment route.} A second route summed Haar-random permutation matrices and looked at the spectrum of the real part. At $n=64$ the standardised free cumulant of order four differed between $k=2$ and $k=8$ by approximately $1.94$, with a pooled standard error of approximately $0.0344$, i.e. a ratio of about $56.34$. This separates $k$ at fixed $n$, but the separation is finite-$n$: the normalisation of the model contributes a factor $-1/(2k)$ to the cumulants, so a difference in $k$ is an arithmetic consequence of the normalisation and not an invariant of the tuple. We therefore record the reading as a heuristic observation about the model, not as evidence for non-isomorphism.

\paragraph{A.5 Gate criteria and the last run.} To decide whether a sequence of estimates could support any conclusion at all, four gate criteria were imposed on an estimator: (Q1) convergence, the relative standard error non-increasing in $N$ with final value at most $0.1$; (Q2) $N$-scaling consistency, the relative change of the estimate from $N=32$ to $N=64$ at most $0.2$; (Q3) normalisation invariance, relative difference at most $0.05$ between normalisation scales at fixed $(k,N)$; (Q4) feasibility, at most one level reported as infeasible. The parameter range was $k\in\{2,3,4\}$, $m\in\{3,4,5\}$ with $m\neq k$, and $N\in\{8,16,32,64\}$, with a fixed deterministic seed per level; $7$ cells were measured. Readings with convergent error were obtained at $k=2,N=8$ (relative standard error $0.0348$) and at $k=3,N=8$ and $k=3,N=16$ ($0.0067$ and $0.0308$). All levels with $k=4$ returned an estimate equal to $1.0$ and carried no information. The cells $k=2$ with $N=16,32,64$ and $k=3$ with $N=32,64$ were rejected by an analytic hit-rate precheck as infeasible. The deterministic gate therefore returned the verdict `insufficient data' with the gate disqualified, the failing reason being that no estimator passed Q1 together with Q2 and Q3; the obstruction is the absence of a qualified estimator, not a shortage of data points.

\paragraph{A.6 What would change the picture.} An estimator satisfying (i)--(iv) of Open problem~\ref{op:estimator} would allow the value $k$ to be checked at finite level; a positive check would not prove Conjecture~\ref{conjInv}, but a failure of stability in $N$ at every feasible level would show that the microstates route cannot be pushed numerically at accessible sizes.

\end{document}
